% Vietnamese translation for OI Public Library
% Source-File: IOI中国国家候选队论文/2004/肖天.ppt
\documentclass[11pt]{article}
\usepackage{oipl-vn}

\newcommand{\Oh}{\mathcal{O}}

\title{Tư tưởng đồ thị phân tầng\\{\large Ghi chú theo slide}}
\author{Xiao Tian}
\date{IOI 2004}

\begin{document}
\maketitle

\origtitle{Layered graph thinking}
\sourcepath{IOI 2004 / Xiao Tian.ppt}

\section{Ý tưởng}

Đồ thị phân tầng tách trạng thái thành nhiều lớp, mỗi lớp ghi một phần lịch
sử hoặc tài nguyên đã dùng. Khi đường đi trong đồ thị gốc không đủ thông tin,
ta nhân bản đỉnh theo trạng thái phụ để biến bài toán thành đường đi ngắn
nhất hoặc quy hoạch động trên đồ thị mới.

\section{Các mốc trên slide}

Slide nhắc các bài như \emph{SuperMan}, \emph{Dad's Tie} và \emph{Rescue}.
Một ví dụ có \(N=M\le 20\), \(P=3\), trạng thái
\[
  G(s_1,s_2,s_3),\qquad s_i\in\{0,1,2\},
\]
với \(G(0,0,0)\) là nguồn và \(G(2,2,2)\) là đích. Vì có \(3^3=27\) lớp cho
phần trạng thái phụ, bài toán nhỏ này minh họa rõ cách nhân bản không gian
trạng thái.

Các độ phức tạp \(\Oh(N^4)\), \(\Oh(N^3)\) và \(\Oh(N^2)\) trên slide cho
thấy mục tiêu của mô hình phân tầng: sau khi thêm trạng thái đúng, ta có thể
áp dụng thuật toán đồ thị chuẩn và tiếp tục giảm bậc bằng quan sát cấu trúc.

\section{Khi nào dùng}

Nên nghĩ tới đồ thị phân tầng khi bài toán có câu như ``được dùng tối đa
\(k\) lần'', ``đã đi qua loại cạnh nào'', ``mang theo trạng thái nào'', hoặc
``hai hành trình diễn ra đồng thời''. Trạng thái phụ phải đủ để quyết định
chuyển tiếp kế tiếp, nhưng càng nhỏ càng tốt để số lớp không bùng nổ.

\end{document}
